
\noindent Large-scale quantum computers have the potential to revolutionize
scientific and engineering computations. Critical to realizing
this potential is the compilation of a quantum
algorithm into a circuit composed of operations natively executable on a quantum computer, and subsequent optimization of this circuit taking into account constraints of the target hardware. These steps must be done for circuits operating on hundreds and thousands of qubits for applications of impact. We refer to this
as \emph{application scale} quantum circuit compilation and optimization. 

Compiled quantum circuits nearly always \emph{approximate} a target unitary prescribed by the quantum algorithm for multiple reasons. 
Firstly, utilizing approximate circuits
in order to reduce circuit depth, especially reductions in expensive
or noisy gates, is often desirable. In the near-term intermediate scale (NISQ) \cite{Preskill2018_NISQ} regime this is a
strategy to increase solution accuracy, while in the
fault-tolerant (FT) regime arbitrary angle rotations \emph{must be}
approximated via a fixed gate-set, and moreover, identifying and
removing insignificant but expensive gates (\eg very small angle
controlled rotations) is a strategy to increase circuit efficiency and
reduce execution time. Secondly, on the theoretical side, exact
circuit compilations are NP-hard to generate with increasing circuit
width \cite{botea2018complexity}, and thus approximations are likely
the best route to scalability.

While there exist many quantum circuit compilation tools \cite{qiskit,bqskit,tket,cirq,qdrift},
none are suitable for application scale compilation, which requires
efficient compilation of circuits with controllable approximations, acting on hundreds and thousands of
qubits. As quantum computing moves from the NISQ
era to early fault-tolerant \cite{PRXQuantum.5.020101} and ultimately full FT devices,
compilation and optimization must evolve to incorporate systematic
approximation mechanisms. Several challenges must be addressed: (i) {\it Quality of optimization:} circuit optimization must balance resource efficiency against approximation error, with tunable cost objectives; (ii) {\it Generality:} approximations must apply to circuits represented at multiple abstraction levels; NISQ-era, continuously parameterized gates and discrete FT gate sets such as Clifford+T; (iii) {\it Soundness;} approximation methods must rest on rigorously proven principles; (iv) {\it Rigorous guarantees:} optimized approximate compilations must come with guarantees on approximation error; and (v) {\it Scalability;} as mentioned above, circuit compilation and optimization tools must address application scale circuits. 

To meet these challenges, we develop and demonstrate an application scale quantum
circuit compilation and optimization workflow. Given an idealized or textbook circuit
expressed in either parameterized or discrete gates, our method
outputs an \emph{ensemble} of approximate compilations tailored to a
target device specification. The method supports arbitrary target error
levels and program execution consists of sampling circuits from the
ensemble and averaging measurement outcomes. This averaging defines a
quantum channel, and we derive rigorous error bounds for this channel
relative to the target unitary.

Our approach combines three main ingredients: (i) partitioning the
input circuit into non-overlapping subcircuits using binning
techniques that group adjacent gates into fixed-size blocks, (ii)
approximate compilation of each partition into an ensemble of
circuits, and (iii) optimization of the ensembles to guarantee circuit
output accuracy. The classical overhead of all steps scales
polynomially with circuit width and depth.

While these ingredients have been individually explored in previous work, their
integration into a complete, practical compilation workflow is our
main contribution. Circuit partitioning as a divide-and-conquer
strategy was previously proposed in Refs.~\cite{wu_qgo,Burt_2024}, but
these works neither analyzed approximation error nor established
output error guarantees.

The idea of ensemble averaging to reduce output error originates with
Campbell~\cite{PhysRevA.95.042306} and
Hastings~\cite{hastings2016turning}, who showed that properly designed
ensembles can suppress worst-case error quadratically
($\epsilon \rightarrow \epsilon^2$). Their methods, however, were
demonstrated only for single-qubit gates and relied on bespoke
compilation techniques (\eg Campbell's ensemble is
specified in terms of the Hamiltonians generating the gates or
circuits \cite{PhysRevA.95.042306}). Following this work, the insight
that averaging over several compilations of a circuit or gate can
reduce error has been applied in several contexts.  It has been
applied to generate randomized single qubit gate compilation
protocols \cite{Kliuchnikov_2023,Akibu_ACM_2024,yoshioka2024errorcrafting},
to define randomized versions of state preparation
circuits \cite{Akibue2024}, to randomize specific computational
primitives \cite{low2021halving}, and specific quantum algorithms,
such as product formula-based quantum
simulation \cite{campbell_qdrift,
Ouyang2020compilation,PRXQuantum.2.040305,huang_2023,PhysRevResearch.6.013224,PhysRevA.109.062431}
and quantum signal processing-based quantum
simulation \cite{martyn2024}.  Prior work by some of the authors
integrated circuit partitioning, synthesis, and
ensembles~\cite{patel_2021}, but lacked application-level scalability,
was primarily empirical, and did not derive worst-case error bounds or
methods for ensemble optimization.

In this work, we integrate theoretical results on error reduction
through circuit ensembles with practical compilation tools, producing
a scalable framework for compiling and optimizing quantum circuit ensembles with
rigorous worst-case output error guarantees.









